\section{Goal Extraction via Reflection}
\seclabel{reflection}
\newcommand\prove[1][]{\UseVerb[#1]{prove}}
\begin{SaveVerbatim}[commandchars=\\\{\}]{prove}
\IdrisFunction{Elab}
\end{SaveVerbatim}

\newcommand\Elab[1][]{\UseVerb[#1]{Elab}}
\begin{SaveVerbatim}[commandchars=\\\{\}]{Elab}
\IdrisType{Elab}
\end{SaveVerbatim}

\newcommand\TC[1][]{\UseVerb[#1]{TC}}
\begin{SaveVerbatim}[commandchars=\\\{\}]{TC}
\IdrisType{TC}
\end{SaveVerbatim}

\newcommand\TTImp[1][]{\UseVerb[#1]{TTImp}}
\begin{SaveVerbatim}[commandchars=\\\{\}]{TTImp}
\IdrisType{TTImp}
\end{SaveVerbatim}

Thus far, our examples have illustrated interaction with \Frexlib{}
using \solve{}. The \solve{} function provides a similar interface to
the simplifiers in (e.g.)  Agda's standard library: it takes the fral
or frex simplifier, the number of free variables and the abstract
syntax of a goal. However, these simplifiers additionally provide
ergonomic goal extraction with Agda's \emph{proof reflection}
mechanism.

Proof reflection is a metaprogramming paradigm, available in
proof assistants and dependently typed programming languages, that
supports bi-directional communication between a language and its
implementation. The language provides
a representation of its terms,
operators that construct, manipulate and destruct term representations, and
primitives \emph{quote} and \emph{unquote} that respectively
\emph{reify} terms into the representation and
\emph{reflect} back encoded terms as ordinary terms.

Given mechanisms for querying unsolved proof obligations, proof
reflection enables the implementation of verified decision procedures
for automatically discharging such obligations without
boilerplate~\cite{boutin:reflection,christiansen-et-al:elab-reflection}.
Coupled with the meta-theoretic properties that
dependently typed implementations of decision procedures can enforce
(e.g.~relative soundness and completeness), reflection-driven interfaces
yield easy-to-use tactics with strong guarantees. It is then natural to
ask: is it possible to construct an interface to \Frexlib{}
that uses proof reflection to avoid the need to explicitly supply the
equation to discharge?

\begin{figure}
  \begin{tabular}{@{}c@{}|c@{}}
  \begin{minipage}[b]{.55\textwidth}
    \MagicExOne{}
    \vspace{-1\baselineskip}
  \end{minipage}
  &
  \begin{minipage}[b]{.44\textwidth}
    \vspace{-1\baselineskip}
    \begin{code}
    \>[0]\AgdaFunction{agdaEx}\AgdaSpace{}%
    \AgdaSymbol{:}\AgdaSpace{}%
    \AgdaSymbol{∀}\AgdaSpace{}%
    \AgdaSymbol{\{}\AgdaBound{x}\AgdaSpace{}%
    \AgdaBound{y}\AgdaSymbol{\}}\AgdaSpace{}%
    \AgdaSymbol{→}\AgdaSpace{}%
    \AgdaSymbol{(}\AgdaNumber{2}\AgdaSpace{}%
    \AgdaOperator{\AgdaPrimitive{+}}\AgdaSpace{}%
    \AgdaBound{x}\AgdaSymbol{)}\AgdaSpace{}%
    \AgdaOperator{\AgdaPrimitive{+}}\AgdaSpace{}%
    \AgdaSymbol{(}\AgdaBound{y}\AgdaSpace{}%
    \AgdaOperator{\AgdaPrimitive{+}}\AgdaSpace{}%
    \AgdaNumber{3}\AgdaSymbol{)}\AgdaSpace{}%
    \AgdaOperator{\AgdaDatatype{≡}}\AgdaSpace{}%
    \AgdaBound{x}\AgdaSpace{}%
    \AgdaOperator{\AgdaPrimitive{+}}\AgdaSpace{}%
    \AgdaSymbol{(}\AgdaBound{y}\AgdaSpace{}%
    \AgdaOperator{\AgdaPrimitive{+}}\AgdaSpace{}%
    \AgdaNumber{5}\AgdaSymbol{)}\<%
    \\
    \>[0]\AgdaFunction{agdaEx}\AgdaSpace{}%
    \AgdaSymbol{=}\AgdaSpace{}%
    \AgdaMacro{fragment}\AgdaSpace{}%
    \AgdaFunction{CSemigroupFrex}\AgdaSpace{}%
    \AgdaFunction{+-csemigroup}\<%
  \end{code}
  \end{minipage}
  \end{tabular}
  \caption{Goal extraction in (a) \idris{} \Frexlib{}'s elaborator reflection script;
    (b) the \textbf{fr}ex \textbf{Ag}da aug\textbf{ment}ation lib.}
  \figlabel{magic}
\end{figure}
As the example code in \figref{magic} illustrates, using proof
reflection to provide such an interface to \Frexlib{} is possible in
both \idris{} and \agda{}.
%
Rather than designing custom reflection-based drivers for individual
simplifiers, we combine proof reflection with \Frexlib{}'s design
philosophy of extensibility and common core reuse and provide a single
generic metaprogram parameterized by a signature and a model of a
presentation.
%
The metaprogram can be instantiated for any algebraic simplifier,
built-in or user-defined.
%
\figref{magic} (a) shows the invocation of the \idris{} elaboration
script \IdrisFunction{frexMagic}, and \figref{magic} (b) shows the
\agda{} proof reflection macro \AgdaMacro{fragment}.
%
Both implementations aim to infer the abstract syntax of the goal
equation based on the expected type.

%% The \idris{} and Agda implementations of \Frexlib{} contain
%% reflection-based drivers with identical interfaces. Each driver
%% receives a frex simplifier and a model of the corresponding
%% presentation.

The drivers have no information about the structure of the algebraic
signature argument ahead of time. \Frexlib{}'s inductive \Term{}
representation means that relevant abstract operator names can be
extracted from the presentation.
%
However, the process of matching goal fragments against the abstract
syntax of the algebraic interpretation is tightly coupled to the
language's reflection primitives. Implementing \Frexlib{} in both
\idris{} and \agda{} allows us to compare differences in behaviour.

The differences between the \idris{} and \agda{} implementations can
be seen by considering the normalisation of arithmetic expressions
such as \ExOne{}.
%
In \idris{}, the reflected syntax passed to the driver represents the
normalized expression \ExTwo{}.
%
As far as the theory of monoids is concerned, \ReflectionSY{} is an atomic
expression and so it is treated as another free variable, distinct
from \ReflectionY{}. The \idris{} driver then incorrectly infers the
invalid equation \ExThree{}, and fails to discharge the goal.
%
In contrast, Agda does not normalise quoted expressions before
reflection, and so the Agda driver successfully finds the equation,
allowing \Frexlib{} to solve the example.
%
Agda's approach is not always superior: there are
similar examples for which the \agda{} driver fails.  The extent to which
the implementation can avoid such pathologies ultimately depends on
the engineering effort available to develop heuristics.

As these problems indicate, this rather naive approach to automation
requires significant developer resources to deal with edge cases or
construct bespoke solutions under simplifying assumptions.
%
In large, mature ecosystems it may be possible to maintain practical
heuristics despite these challenges.
%
However, there might be better mechanisms for specifying
algebraic contexts, allowing the solver to extract the required
information automatically; we suggest such directions in
\secref{conclusion}.
